% Copyright 2026 icekylinx and GPT-6 Astra Pro. Apache-2.0.
% Curated from sat_research_checkpoint.zip: paper/integral_schur_extension.tex
% Complete integral crossing theorem and proof retained; unrelated downstream corollaries and duplicate streaming implementation omitted.
\subsection{An integral Schur basis removes the factorial normalization}
\label{sec:integral-schur}

The factorial $h!$ in the source's monomial-basis construction is
not intrinsic to its rank bound. A change to an integral Schur
basis removes it. Integral forms of Heisenberg and lattice vertex
operators are established theory~\cite{huangjing2022}; the proof
below gives the specific finite coefficient statement needed here.
The required source operator identities are independently proved
in Lemmas~\ref{lem:elementary-crossing-operator}
and~\ref{lem:elementary-crossing-duality}. This algebraic result
does not depend on the source SAT algorithm.

\begin{theorem}[Integral crossing pieces]
\label{thm:integral-crossing}
For the modes defined in Section~\ref{sec:elementary-crossing-operator}
and all integers $h,E\ge1$, there are
integer crossing pieces indexed by $-E-h\le s\le E+h$ satisfying
\[
 \sum_s A_s(x,y,u,v)=[x=v][y=u],
\]
where $x,v$ range over the partitions of $h$ and $y,u$ over those
of $E$. Both factorizations of every piece have integer entries,
and their ranks are at most
\[
 R_s^{\rm row}=\sum_{i=0}^{2h}p(E+s-h+i),\qquad
 R_s^{\rm col}=\sum_{i=0}^{2h}p(E+h-s-i).
\]
Here $p(j)=0$ for $j<0$. Thus the normalizing scalar is $1$, and
the identities and rank bounds specialize to every commutative
ring. The integers are explicitly constructible with bit lengths
polynomial in $h+E$.
\end{theorem}

\begin{proof}
Identify the source variables $t_j$ with power-sum symmetric
functions $p_j$. Put
\[
 H(u)=\sum_{n\ge0}h_nu^n
     =\exp\!\left(\sum_{n\ge1}p_nu^n/n\right),
 \qquad E(u)=H(-u)^{-1}.
\]
The integral lattice is the symmetric-function ring
$\Lambda_{\mathbb Z}=\mathbb Z[h_1,h_2,\ldots]$. Products
$h_\lambda$ and Schur functions $s_\lambda$ are integral bases.
The Hall pairing has
\[
 \langle p_\lambda,p_\mu\rangle
   =z_\lambda[\lambda=\mu],\qquad
 \langle s_\lambda,s_\mu\rangle=[\lambda=\mu].
\]

First, every source mode of an element of $\Lambda_{\mathbb Z}$
preserves $\Lambda_{\mathbb Z}$. For a homogeneous state $a$ of
weight $w$, write $Y(a,z)=\sum_r a[r]z^{r-w}$ and extend the
state variable linearly. The source's normal ordering places all
multiplications before all differentiations. Direct summation of
its binomial coefficients gives
\[
 Y(H(u),z)f
 =\frac{H(z+u)}{H(z)}\,
 f\left[p_k\mapsto p_k+z^{-k}-(z+u)^{-k}\right].
\]
Indeed the creation coefficient of $p_l$ is
\[
 \sum_{n\ge1}\frac{u^n}{n}\binom{l-1}{n-1}z^{l-n}
 =\frac{(z+u)^l-z^l}{l},
\]
and the coefficient of $\partial_{p_k}$ in the annihilation
exponential is
\[
 k\sum_{n\ge1}\frac{u^n}{n}\binom{-k-1}{n-1}z^{-k-n}
 =z^{-k}-(z+u)^{-k}.
\]
More generally, normal ordering once across all factors yields
\begin{equation}
\label{eq:integral-mode-generating}
 Y\!\left(\prod_{j=1}^qH(u_j),z\right)f
 =\frac{\prod_{j=1}^qH(z+u_j)}{H(z)^q}\,
 f\left[p_k\mapsto p_k+qz^{-k}
                       -\sum_{j=1}^q(z+u_j)^{-k}\right].
\end{equation}
The substitution adds $q$ copies of the letter $z^{-1}$ and
subtracts the letters $(z+u_j)^{-1}$ in the symmetric-function
alphabet. It is integral: on the complete-function generating
series it multiplies by
\[
 \frac{\prod_j(1-(z+u_j)^{-1}t)}{(1-z^{-1}t)^q}.
\]
The prefactor in~\eqref{eq:integral-mode-generating} is integral
because $H(z)^{-1}=E(-z)$. Fixing the multidegree in the $u_j$ and
the weight of $f$ gives a lower Laurent bound in $z$; every
requested coefficient is a finite sum. Extracting a monomial in
the $u_j$ therefore proves mode integrality for $h_\lambda$, and
hence for every integral symmetric function. In particular, all
mode matrices in the Schur basis are integer matrices.

Next define the source's triangular transform
\[
 T(p_n)=(-1)^{n-1}\sum_{j=1}^n\binom njp_j.
\]
It obeys
\[
 T(H(u))=E\!\left(\frac{u}{1+u}\right),\qquad
 T(h_n)=\sum_{j=1}^n(-1)^{n-j}\binom{n-1}{j-1}e_j.
\]
Thus $T$ preserves the integral weight filtration, including
integrality of each homogeneous component. The finite Vandermonde
calculation in Lemma~\ref{lem:elementary-crossing-duality} gives
\[
 (T p_\lambda)[-h]p_\mu=z_\lambda[\lambda=\mu]
 \quad (|\lambda|=|\mu|=h).
\]
Extending over $\mathbb Q$ identifies this pairing with the Hall
pairing. Consequently
\[
 a_\lambda=T(s_\lambda),\quad b_\nu=s_\nu
 \quad\Longrightarrow\quad
 a_\lambda[-h]b_\nu=[\lambda=\nu].
\]

Use these $a_\lambda,b_\nu$ in the source's $Q_s$ and take matrix
entries in the Schur basis of $V_E$. The state
\[
 a_\lambda^{\langle i\rangle}
 =\sum_{w=0}^h\binom{h+w}{i}(a_\lambda)_w
\]
is integral componentwise. The mode-integrality argument therefore
makes every entry of both factorizations an integer.
Lemma~\ref{lem:elementary-crossing-operator} gives
the stated crossing identity with scalar $1$. The two intermediate
weights are $E+s-h+i$ and $E+h-s-i$, so the dimensions give exactly
the asserted ranks. Standard finite symmetric-function basis
conversion, or the generating identities above, constructs all
entries; factorial denominators in an intermediate power-sum
calculation cancel in the final integer coordinates, and all
intermediate weights are $O(h+E)$, giving polynomial bit height.
\end{proof}
